We built a system that chose stocks, set their share counts by maximizing a utility under limits on weight,
volatility and turnover, and re-planned the holdings with a model-predictive controller, and we tested it on $373$
United States equities over $420$ trading days of 2025--2026 and on simulated markets. On real prices, the choice of
stocks moved terminal wealth more than the choice among seven weighting rules: one $33$-stock selection ended between
$1.09$ and $1.23$ times initial wealth across the weighting rules, while different selections ended between $1.03$
and $1.70$, the highest being a $13$-stock list the author picked in advance. The share of the controller's re-plans that came
before the $21$-day horizon, because the portfolio value left the $95\%$ band of its forecast or fell $8\%$ below its
peak since the last drawdown stop, ranged from a fifth to three fifths across combinations of
selection and weighting rule, highest for stocks selected by their score, the expected growth rate from a
single-index fit divided by a power of the stock's market beta that depended on the market's recent direction.
Compared with re-planning every $21$ trading days alone, these early re-plans made the median worst drawdown smaller
in $55$ of the $70$ combinations of the ten selection rules and seven weighting rules without the market hedge and
in $61$ with it, and raised trading cost in all of them, but raised median terminal wealth in $38$ and $26$ of the
$70$. Among $10{,}000$ simulated price paths (synthetic futures), we kept two sets of $200$ with the
deepest drawdowns of simulated SPY, one from futures whose SPY went through a burst of extreme daily moves and one
from futures without such a burst. In both sets, the median terminal wealth of every combination ended above those
of simulated SPY bought and held and of the equal-weight universe. Random selections did so as well, and so did
every combination with the band exit and the drawdown stop switched off. The equal-weight universe also ended above
simulated SPY, so part of the margin over SPY came from the simulated market and the rest from parts of the system
that every combination shared other than those two triggers, such as the volatility cap, whose separate effect the
study did not measure. In these sets the early re-plans again made the median worst drawdown smaller, in $65$ to
$70$ of the $70$ combinations for each set and hedge setting, without consistently raising median terminal wealth. In $600$ futures drawn at random from the $10{,}000$ (typical futures),
the $13$ stocks with the highest scores ended above $13$ random stocks in $62$ to $71\%$ of futures across the
weighting rules, but the generator kept each stock's 2014--2024 trend constant, and the score's trend term estimates
that trend, whereas in the real market those trends did not carry over into 2025--2026, so
part of this advantage came from the generator's design and may not carry over to real prices. A market hedge, which at
each re-plan capped the portfolio's beta more tightly when the market was weaker and added stocks with zero or negative beta, reduced
terminal wealth on real prices and added to it in the set without bursts, whose declines were the most sustained. Deciding which selection rule is
better, and when such a hedge pays, will need more historical windows and generators whose trends persist no more
than real ones do.
